\paragraph{Inventory and evidence strength.}
The historical inventory contains 17 experimental series and 171 evaluation records, including alternative grader views and explicit missing results. It is not a count of 171 independent training runs. Below, percentages are reported separately for each benchmark. A dash means that no completed result was recovered, not zero accuracy. In the n=8 tables, cells contain Avg@8 / Pass@8. The early n=1 and 100-question pilot tables retain their original metric definitions. Regrading the same generations does not add samples, and re-evaluating an old checkpoint does not add a training seed.

The strongest retained evidence consists of original evaluation summaries with their acceptance records. Other tables come from finalized paired artifacts or curated summaries. The earliest block-sum and block-advantage results are available only as rounded values in historical reports; we retain that limitation rather than reconstructing exact counts from rounded percentages. Source paths and identifying infrastructure information remain outside the anonymous manuscript. A complete internal provenance map records which fields support each value.

\paragraph{H01--H05: clean-room exploration.}
The earliest student and teacher were the official post-trained \texttt{Qwen3-0.6B} and \texttt{Qwen3-4B}, not their Base variants and not the later GRPO teacher. The 3,400-prompt training pool combined 2,000 GSM8K training examples with 200 training examples from each of MATH's seven subjects. Small weighting pilots, outcome weighting, reference anchoring, naive block sums, and sum/mean/mixed variants were explored. Full evaluation meant GSM8K's 1,319 test examples and 500 MATH examples, n=1, with a 256-token output cap, temperature 0.7 and top-$p=0.95$. The 100-example pilots are not full evaluations and their sample-count parameter is not recoverable from those summaries.

These runs use a different loss family from the main study. The clean-room implementation applies detached weights and advantages to summed current log probabilities without a PPO ratio, clamps its raw signal, and discards incomplete blocks. Its mixed mode interpolates token-local and block-mean signals; this is not the later block-sum/block-mean interpolation option. Exact settings for some named router or anchor variants are incomplete. Consequently, the old pilot tables document the research trajectory and negative outcomes, not replications of the current loss or implementations of similarly named external papers. Outcome weighting did not prevent endpoint degradation: the two reweighted 200-step arms attain primary-grader MATH500 accuracies of 16.60\% and 16.40\%, below the recorded 20.20\% initial reference; the standard arm attains 14.80\%. Missing intermediate reweighted evaluations prevent locating the onset of that degradation. The standard arm's Step50 entry in an older composite chart came from an earlier pilot and is not connected here as if it were the same training run.

\paragraph{H06--H09: size sweep, collapse and Qwen 1.7B.}
The DAPO-era Qwen3-1.7B-Base student uses the retained 4B GRPO teacher. The original block-size sweep mixed hardware contexts and did not explicitly fix evaluation seeds in every older invocation. Token, Block3 and Block5 used A800 hardware; Block10 used A100 hardware. The retained Step200 Avg@8 / Pass@8 table for $k\in\{1,3,5,10\}$ appears in Appendix~\ref{history:qwen17_sweep}. Block3 exceeds Block5 on every Avg@8 metric and on three Pass@8 metrics, tying AMC23 coverage. Its comparison to Token is mixed, so the sweep motivates retaining a short block rather than claiming universal superiority. The Block5 scores survive as a curated summary rather than recovered original evaluation files. The first Block5 attempt stopped at Step42 with a vLLM OOM; the successful restart lowered the rollout-memory fraction to 0.6. These resource changes must not be hidden when interpreting the sweep.

An earlier clean-room pilot compared $k\in\{1,2,3\}$ at 50 updates (Appendix~\ref{history:early_blocksum50}). Its recorded GSM8K/MATH500 accuracies were 45.94/20.00 for Token, 45.41/19.40 for Block2 and 46.85/20.20 for Block3. It used official 0.6B/4B post-trained models and a non-PPO loss, so this is preliminary motivation, not a $k=2$ result for the later 1.7B PPO implementation. No completed $k=2$ arm exists in the retained 200-step sweep.

The later same-host Token/Block3 comparison provides the cleaner historical method comparison. At Step200, its MATH500 Avg@8 changes from 45.85\% to 54.75\%, AIME24 from 5.42\% to 8.75\%, AIME25 from 2.50\% to 5.42\%, and AMC23 from 21.84\% to 27.26\%. Intermediate checkpoints are not uniformly favorable. Six later checkpoint evaluations regenerate outputs for the same Token/Block3 weights; these are evaluation sensitivity records, not six new training runs.

Both 1.7B arms use ChatML and request \texttt{<think>...</think>} during training without explicitly disabling thinking in the template kwargs. Evaluation uses ChatML with thinking disabled and no instruction requiring those tags, retaining the already-closed empty thinking-control prefix. The old collector resets the per-request training seed to 21, so nominal eight-response groups may contain duplicates. The original full training-token trajectories were not retained, preventing retrospective measurement of their exact duplicate rate. These shared historical choices do not become the current completion protocol merely because the dataset and global seed match.

The Block10 diagnostic repetitions both use training seed 21, one on A800 hardware and one on A100 hardware. Both record zero scores on all four benchmarks at Step200 while available optimizer diagnostics remain finite. Their grader is the legacy built-in VERL path, not the pinned external path used for the current comparison. In particular, built-in AMC23 zero scores have a known compatibility caveat and cannot independently establish absence of AMC capability; we retain them as historical outputs, not validated capability measurements. The severe generation degradation and failures on the other tasks remain relevant observations. These are two hardware-context repetitions, not a multiple-training-seed confidence estimate. Neither the scalar summaries nor these endpoints prove a unique cause such as gradient explosion, teacher unreliability, or reverse-KL mode collapse.

Table~\ref{tab:block10health} makes the limitation of a set-overlap diagnostic concrete: the two runs retain top-16 overlap ratios near 0.8 at Step200, while the shared tokens carry only about 0.17--0.24 of the full probability mass. High rank-set overlap need not imply concentrated agreement. The largest finite pre-clipping gradient norms occur at different steps (191 on A800 and 1 on A100), so these summaries do not support one reproducible temporal ordering of collapse.
\begin{table}[ht]
\centering\small
\caption{Retained Block10 diagnostic summaries. Entropies (nats), overlap ratio and intersection probability masses are measured at Step200; the last column is the maximum pre-clipping gradient norm over each run. Values come from the archived diagnostic report rather than a fresh training replication.}
\label{tab:block10health}
\begin{tabular}{@{}lrrrrrr@{}}
\toprule
Hardware & $H_S$ & $H_T$ & Overlap & $M_S$ & $M_T$ & Max. grad \\
\midrule
A800 & 7.146 & 6.745 & 0.837 & 0.226 & 0.237 & 867.495 \\
A100 & 7.576 & 6.962 & 0.793 & 0.166 & 0.186 & 865.207 \\
\bottomrule\end{tabular}

\end{table}

\paragraph{H10--H11: another distilled pair and alternative windows.}
The DeepSeek-R1-Distill-Qwen-1.5B/JustRL pair is a separate released-model comparison and still has a Qwen-derived backbone; it is not a new non-Qwen architecture. Its Step200 Avg@8 differences favor Block3 by 0.05, 4.17, 1.67 and 0.45 percentage points on MATH500, AIME24, AIME25 and AMC23, respectively. However, Step200 Pass@8 decreases on MATH500 and AIME24; Step50 AIME24/AMC23 and Step100 AIME25 also contain regressions. The archived protocol classifies this as a failed directional replication: its original criterion required positive equally weighted task-mean differences for both Avg@8 and Pass@8. We preserve that historical decision without using an aggregate score in the present tables. Positive Avg@8 endpoints should neither erase the failed criterion nor be described as uniformly negative. One training seed does not establish general transfer; all task/checkpoint views remain available.

Random3 and Sliding3 were trained and evaluated as their own pair at Steps 50, 100 and 200. Random3 chooses one of three partition offsets per update using a saved PCG64 RNG (seed 910021); Sliding3 averages the three complete offset-specific losses, after each phase's PPO clipping. Both retain partial edge windows. Their length-weighted phase-zero loss is algebraically equivalent to historical Block3 when valid responses start at position zero with right padding and the $10^{-8}$ denominator stabilizer is ignored. It would be misleading to call this a fundamentally different normalization. The implementation restarts coordinates at mask-separated valid regions and changes floating-point summation, so arbitrary masks need not give identical partitions. Each still performs one optimizer update per batch. A late advantage of Sliding3 over Random3 does not establish superiority over a newly matched fixed Block3 or Token control. Both alternatives can deteriorate during training. The archived grader recovery changes evaluation plumbing, not the frozen training run. Corrected external views support method comparisons; older built-in views remain explicitly separate, including the AMC23 grading caveat above.

\paragraph{H12--H15: small Base and instruction-tuned students.}
The old Qwen3-0.6B-Base attempt trained Token OPD but its Step50 evaluation was stopped during truncation diagnosis; no corresponding Block3 training is claimed. A later paired experiment explicitly disables thinking and retains ChatML, with both arms using the same repaired non-thinking protocol. This repair also changes the EOS stopping set and masks according to actual generated lengths; it is not a prompt-only intervention relative to the older run. It produces a strong negative transfer result for Block3: its Step200 MATH500 Avg@8 is 0.57\%, versus 42.88\% for Token. The 1.7B result is therefore not a universal argument for using blocks in smaller students.

The Llama pair is \texttt{Llama-3.2-1B-Instruct} and \texttt{Llama-3.2-3B-Instruct}, obtained as ModelScope mirrors of the same Meta release \citep{llama32}. It is not a 3.1 teacher paired with a 3.2 student, and neither checkpoint is a Base model. In the first native-prompt attempt, initial-model evaluation completed but Token training was deliberately stopped at Step146. A new pair then reproduced the historical 1.7B training/evaluation instruction difference while retaining Llama-native role boundaries. Token and Block3 complete 200 steps in this setting, with Block3 substantially worse on MATH500 and AMC23. The thinking flag is not a native Llama mode switch; an instruction asking for think tags is separate from the tokenizer's role format.

Both the 0.6B and Llama training series retain the old per-request fixed seed 21. Archived groups show that the eight nominal training responses to one prompt can be duplicates, so 6,400 retained trajectories do not imply 6,400 independent generations. The current completion comparison repairs this request-level seed rule. This training-sampling limitation is distinct from the eight independently seeded evaluation responses. It applies to interpreting effective exploration, not to retroactively changing any recorded benchmark score.

\paragraph{H16--H17: pre-update repetition and the completion restart.}
The 4B Base-model ChatML attempt was stopped after four updates. Repetition and length-cap hits were already present in its initial pre-update outputs, so those initial failures cannot be caused by its later Block3 gradients. Bounded prompt diagnostics examined native input token IDs, thinking instructions, EOS handling and request-seed reuse. The subsequent completion protocol simultaneously removes ChatML and uses independently derived per-request seeds. It is an engineering repair with matched new Token/Block3 arms, not a single-factor causal ablation of role tokens.

A first formal restart failed before its first update with a temporary-directory cleanup error on network storage. The retry moved compiler/temporary caches to executable shared memory while retaining the frozen training code and model/data contract. The failed attempt and its first trajectories remain retained. CPU/GPU acceptance probes that save and resume Steps1--2 are recorded as technical acceptance, not as extra benchmark experiments. The current pair's final results appear in Appendix~\ref{app:current}, replacing the pre-evaluation snapshot in this historical inventory.

\paragraph{Archive inclusion.}
The tables report completed evaluations with retained result artifacts; diagnostic probes and interrupted attempts are identified separately. Later checkpoint sampling is recorded as a new evaluation, not as recovery of the original training trajectories. The interpretation of the integrated update and directions for further comparisons are discussed in Appendix~\ref{sec:limits}.
